% ============================================================
%  TurkLang 2026 — Speech notes for the 14-minute talk
%  Paper: A Gold-Standard Dependency Treebank for the Uzbek Language
%  Authors: S. Matlatipov, M. Kurbanova
%
%  Optimised for printing on a black & white printer:
%    - 12pt serif body, 1.4 line spacing
%    - bold cues only (no colours), clear slide separators
%    - generous left margin for handwritten annotations
%    - timing budget printed next to each slide
%
%  Build with:  pdflatex speech_notes.tex
% ============================================================
\documentclass[12pt,a4paper]{article}

\usepackage[utf8]{inputenc}
\usepackage[T1]{fontenc}
\usepackage{lmodern}
\usepackage[margin=2.2cm,left=3.2cm,right=2.2cm]{geometry}
\usepackage{setspace}
\usepackage{enumitem}
\usepackage{titlesec}
\usepackage{fancyhdr}
\usepackage{xcolor}      % only for \textbf weighting; everything renders in black
\usepackage{microtype}
\usepackage{parskip}

\setstretch{1.40}
\setlist[itemize]{leftmargin=1.3em,topsep=2pt,itemsep=2pt}

% --- Slide section style: big, easy to spot when flipping through pages ---
\titleformat{\section}
   {\normalfont\Large\bfseries\filright}
   {Slide \thesection.}{0.6em}{}
\titlespacing*{\section}{0pt}{1.2em}{0.4em}

% --- Header & footer with running timing ---
\pagestyle{fancy}
\fancyhf{}
\fancyhead[L]{\small TurkLang 2026 — Speech notes (14 min)}
\fancyhead[R]{\small Matlatipov \& Kurbanova}
\fancyfoot[C]{\thepage}
\renewcommand{\headrulewidth}{0.4pt}

% Small helper: timing tag at the start of a slide section
\newcommand{\timing}[2]{%
   \par\noindent\textbf{\small [#1\,$\to$\,#2\,min]}\par\vspace{0.2em}}

% Highlight cue: a phrase to read with emphasis (rendered as bold only — print-safe)
\newcommand{\say}[1]{\textbf{#1}}

% Pause / breathe marker
\newcommand{\pause}{\par\vspace{0.4em}\textbf{[pause]}\par}

\begin{document}

% ------------------------------------------------------------
%  COVER PAGE
% ------------------------------------------------------------
\begin{center}
   {\LARGE\bfseries Speech notes — 14 minutes}\\[0.6em]
   {\large A Gold-Standard Dependency Treebank\\ for the Uzbek Language}\\[0.4em]
   {\normalsize XIV International Conference on Computer Processing of Turkic Languages}\\
   {\normalsize \textbf{TurkLang 2026}}\\[1.2em]
   \textit{Sanatbek Matlatipov \quad \& \quad Mukhabbat Kurbanova}\\
   National University of Uzbekistan, Tashkent
\end{center}

\vspace{1em}
\noindent\textbf{How to use these notes.}
Each section corresponds to one slide. The bracketed times
\textbf{[start\,$\to$\,end]} are the cumulative running time. Phrases in
\textbf{bold} are the ones to emphasise; \textbf{[pause]} marks a short
breath. Plain text is the spoken script — read naturally, not word-for-word.

\vspace{0.5em}
\noindent\textbf{Total budget:} 14 minutes of speech + $\approx$1 minute buffer
for transitions and applause. Aim for roughly 0:50--1:10 per slide.

\vfill\hrule\vspace{0.5em}
\noindent\textit{Print double-sided on A4. Left margin is wide for handwritten cues.}

\newpage

% ------------------------------------------------------------
%  SLIDE 1 — Title
% ------------------------------------------------------------
\section*{Slide 1 — Title}
\timing{0:00}{0:30}

Good morning / afternoon, and thank you to the TurkLang organisers for the
opportunity to present.

My name is \say{Sanatbek Matlatipov}, from the National University of
Uzbekistan, and this is joint work with \say{Mukhabbat Kurbanova}. Today I will
talk about \say{a new gold-standard dependency treebank for Uzbek}.

\pause

The talk will focus on the \say{resource itself and how we evaluate it}, rather
than on parser internals — those are well documented in the paper.

% ------------------------------------------------------------
%  SLIDE 2 — Outline
% ------------------------------------------------------------
\section*{Slide 2 — Outline}
\timing{0:30}{1:00}

I will walk you through the talk in three parts:

\begin{itemize}
   \item First, \textbf{why} a new Uzbek treebank was needed.
   \item Second, \textbf{how} we collected, annotated, and quality-controlled the data.
   \item Third, \textbf{what we learn} from a comparative evaluation with three parsers.
\end{itemize}

\noindent
I will close with conclusions and a short outlook.

% ------------------------------------------------------------
%  SLIDE 3 — Motivation
% ------------------------------------------------------------
\section*{Slide 3 — Why a new Uzbek treebank?}
\timing{1:00}{2:00}

Uzbek is a \say{Turkic, agglutinative, SOV language} with roughly
35 million speakers. Its rich nominal and verbal morphology means that
\say{a single word can carry many grammatical features at once} — case,
possession, tense, mood, person, number.

Until now the Universal Dependencies community had \say{exactly one}
Uzbek treebank: Uzbek-UT, released in 2021 by Akhundjanova and Talamo. It
contains 500 sentences from news and fiction, and was an important first
step.

\pause

Our research question is simple: \say{would a larger, multi-genre,
double-annotated gold treebank meaningfully change parser performance for
Uzbek?}

% ------------------------------------------------------------
%  SLIDE 4 — Related work
% ------------------------------------------------------------
\section*{Slide 4 — Where this fits in Uzbek NLP}
\timing{2:00}{2:45}

Uzbek NLP is a young but active field. On the morphology side, there is
the Prolog analyser from 2009, the rule-based UzbekTagger by
Sharipov and colleagues, and more recent analysers by Salaev and
Abdurakhmonova. On the downstream side, we have sentiment and aspect-based
sentiment analysis work.

\say{What was missing} was an \say{independent, multi-domain, double-annotated}
gold treebank with published inter-annotator agreement. That is the gap we
address.

% ------------------------------------------------------------
%  SLIDE 5 — Contributions
% ------------------------------------------------------------
\section*{Slide 5 — Contributions}
\timing{2:45}{3:30}

In one sentence: we deliver a new gold treebank, with measured
quality, and a head-to-head evaluation of three popular parsers on it.

Concretely:

\begin{itemize}
   \item \textbf{681 sentences, roughly 7\,950 tokens, 4\,800 unique word forms}.
   \item Annotated by \textbf{six annotators} — four linguists and two NLP engineers — fully \textbf{double-annotated} in INCEpTION, with adjudication.
   \item We report inter-annotator agreement using both \textbf{Cohen's kappa} and \textbf{Krippendorff's alpha}.
   \item We benchmark \textbf{UDPipe, Stanza, and spaCy} under an identical protocol.
   \item And the treebank is released publicly in CoNLL-U format.
\end{itemize}

% ------------------------------------------------------------
%  SLIDE 6 — Data sources
% ------------------------------------------------------------
\section*{Slide 6 — Data sources}
\timing{3:30}{4:20}

The data comes from three Latin-script sources:

\begin{itemize}
   \item the contemporary fiction novel \textit{Maqar},
   \item the educational / literary work \textit{Kun shundan boshlanadi}, and
   \item a collection of publicly available Uzbek \textbf{fairy tales}.
\end{itemize}

Sentences are short on average — about \textbf{12 tokens} — which is realistic
for literary Uzbek. The final split is \textbf{481 training sentences and 200
test sentences}, roughly a 70/30 division.

% ------------------------------------------------------------
%  SLIDE 7 — Annotation workflow
% ------------------------------------------------------------
\section*{Slide 7 — Annotation workflow}
\timing{4:20}{5:30}

The annotation pipeline had five stages:

\begin{itemize}
   \item Pre-tokenisation following UD version 2 conventions.
   \item A \textbf{pilot training round} on a shared seed set, so that all annotators converged on the guidelines.
   \item \textbf{Double independent annotation} of every sentence in INCEpTION.
   \item \textbf{Disagreement adjudication}: majority vote first, escalated to a senior linguist when needed.
   \item Export to CoNLL-U.
\end{itemize}

\noindent
We made a few \say{Uzbek-specific decisions} worth flagging: case suffixes are
treated as features rather than separate tokens; \say{postpositions are kept as
separate tokens} and attached with the \texttt{case} relation; and we follow
the UD convention of using \say{no explicit copula node} when Uzbek would
normally drop it.

% ------------------------------------------------------------
%  SLIDE 8 — Annotated example
% ------------------------------------------------------------
\section*{Slide 8 — A real annotated sentence}
\timing{5:30}{7:00}

Let me show you one real example from the treebank — sentence
\texttt{s2}. The English translation is:

\begin{quote}
\textit{`Until now, in studies on dialects, the more differences from the
literary language are shown, the more positively they are evaluated.'}
\end{quote}

\noindent
Three things to notice — these are the \textbf{coloured edges} on the slide:

\begin{itemize}
   \item The postposition \say{\textit{bo'yicha}} — meaning \textit{about / on} —
         is a \textbf{separate token} and attaches to \textit{shevalar} with the
         \texttt{case} relation. This is the recurrent Uzbek
         postposition pattern.
   \item The locative suffix \textit{-da} on \textit{bizda} and \textit{tadqiqotlarda}
         drives two \texttt{obl} (oblique) attachments to the verb.
   \item And the conditional clause headed by \say{\textit{ko'rsatilsa}} attaches
         to the main verb \textit{baholanib} through an \texttt{advcl} edge —
         this is a classic \textit{the more X, the more Y} construction in Uzbek.
\end{itemize}

\pause

I show this example because \textbf{these three relations also drive the
errors we see later} in the evaluation section.

% ------------------------------------------------------------
%  SLIDE 9 — Inter-annotator agreement
% ------------------------------------------------------------
\section*{Slide 9 — Inter-annotator agreement}
\timing{7:00}{8:00}

Quality control. We computed agreement across all six annotators on the
doubly annotated pool.

\begin{itemize}
   \item \textbf{Lemmatisation:} Cohen's kappa of \textbf{95.2\%}, Krippendorff's alpha of 94.6\%.
   \item \textbf{POS tagging (UPOS):} 93.7\% and 94.1\%.
   \item \textbf{Morphological features:} 90.8\% and 89.9\%.
\end{itemize}

\noindent
Two short comments: first, \textbf{lemma and POS are above 93\%} on both
metrics, which is strong agreement on the closed-class layers. Second,
\textbf{morphology is slightly lower}, and this is expected — the
ambiguity in Uzbek derivational suffixes is real, and even trained
linguists sometimes disagree.

% ------------------------------------------------------------
%  SLIDE 10 — Treebank statistics
% ------------------------------------------------------------
\section*{Slide 10 — Treebank statistics}
\timing{8:00}{8:45}

How does the new resource compare to Uzbek-UT in size and coverage?

\begin{itemize}
   \item \textbf{681} sentences versus 500 — about a 36\% increase.
   \item \textbf{$\approx$ 7\,950} tokens versus 5\,850.
   \item \textbf{4\,800} unique surface forms versus 3\,523.
   \item All \textbf{17 UD parts of speech} are covered.
   \item \textbf{45 morphological feature values} (up from 42).
   \item \textbf{34 dependency labels} (up from 32) — we add
         \texttt{discourse}, \texttt{parataxis}, \texttt{vocative}, and
         \texttt{ccomp}, all driven by literary dialogue.
\end{itemize}

% ------------------------------------------------------------
%  SLIDE 11 — Linguistic profile
% ------------------------------------------------------------
\section*{Slide 11 — Linguistic profile}
\timing{8:45}{9:30}

A quick linguistic profile so you can compare the two treebanks at a
glance:

\begin{itemize}
   \item NOUN dominates at around \textbf{30\%}, VERB at 20\%, ADJ at 15\%.
   \item Proper nouns are higher than in Uzbek-UT — \textbf{8\% versus 5\%} — because of named characters in the fairy tales.
   \item Numerals are also slightly elevated, reflecting the educational source.
   \item About \textbf{52\% of tokens carry morphological features}, consistent with the agglutinative profile.
   \item The most frequent relations are \texttt{nsubj}, \texttt{case},
         \texttt{obl}, and \texttt{obj} — again \say{case-marked obliques are everywhere}.
\end{itemize}

% ------------------------------------------------------------
%  SLIDE 12 — Experimental setup
% ------------------------------------------------------------
\section*{Slide 12 — Experimental setup}
\timing{9:30}{10:15}

For the evaluation we deliberately keep things simple and reproducible.
We pick three widely used UD parsers and \textbf{treat them as black boxes}:

\begin{itemize}
   \item \textbf{UDPipe} — a classical pipeline baseline.
   \item \textbf{Stanza} — a neural pipeline with a character-level CNN.
   \item \textbf{spaCy} — a transition-based parser with tok2vec embeddings.
\end{itemize}

\noindent
We train each parser on \textbf{both treebanks separately} using the same
protocol: gold tokenisation, default hyperparameters, and the same
training / test split. The metrics are the standard UD ones:
\textbf{LAS, UAS, and UPOS accuracy}.

% ------------------------------------------------------------
%  SLIDE 13 — Results
% ------------------------------------------------------------
\section*{Slide 13 — Parsing results across both treebanks}
\timing{10:15}{11:30}

This is the main results table. \say{Same protocol, same parsers — only the
training treebank changes.} I want to read it carefully so we are honest
about what the numbers mean.

\begin{itemize}
   \item On \textbf{Uzbek-UT}, the parsers reach \textbf{52 / 60 / 62 LAS} for UDPipe, spaCy and Stanza. That is a \textbf{strong baseline} — every parser already crosses 50 LAS.
   \item On our \textbf{literary / fairy-tale treebank}, the same parsers reach \textbf{55 / 61 / 66 LAS}.
   \item \textbf{Stanza is the best on both treebanks}, which is expected — the character-level features are a natural fit for Uzbek morphology.
\end{itemize}

\pause

\say{I would not call this a competition between the two resources.} They
cover \textbf{different domains}, and they have \textbf{different annotation
protocols} — ours adds double annotation and published IAA. The two
treebanks are \textbf{complementary, not competitive}, and the natural next
step is to combine them.

% ------------------------------------------------------------
%  SLIDE 13b — Training curve
% ------------------------------------------------------------
\section*{Slide 13b — How the Stanza parser learns Uzbek}
\timing{11:30}{12:15}

This plot shows the Stanza parser's \textbf{dev UAS and LAS curves} when
trained on the two treebanks combined. Two observations:

\begin{itemize}
   \item UAS climbs to around \textbf{72} and LAS to around \textbf{63} \say{within the first thousand training steps}.
   \item Both curves \textbf{plateau early} — the model is \say{not under-trained, it is data-limited}.
\end{itemize}

\noindent
This is the clearest evidence we have that the bottleneck for Uzbek
parsing is not the architecture — it is \textbf{annotated data}.

% ------------------------------------------------------------
%  SLIDE 14 — Error analysis
% ------------------------------------------------------------
\section*{Slide 14 — Error analysis}
\timing{12:15}{13:00}

Where do the parsers still fail? Three recurring confusions, common to
all three parsers:

\begin{itemize}
   \item \textbf{\texttt{advcl} clause boundaries} — exactly the kind of conditional clause we saw in the example sentence.
   \item \textbf{Coordination} attachment in long nominal phrases.
   \item Confusion between \texttt{nmod} and \texttt{obl} for case-marked nominals.
\end{itemize}

\noindent
The driver behind all of these is Uzbek's \textbf{flexible word order} on
top of a basic SOV pattern. The natural direction forward is
\textbf{contextual embeddings} — multilingual BERT, XLM-R, or Turkic-specific
models like TURNA.

% ------------------------------------------------------------
%  SLIDE 15 — Conclusion
% ------------------------------------------------------------
\section*{Slide 15 — Conclusion}
\timing{13:00}{13:45}

To summarise:

\begin{itemize}
   \item We deliver a new \textbf{gold-standard Uzbek UD treebank} — 681 sentences, eight thousand tokens, fully double-annotated.
   \item Inter-annotator agreement is \textbf{above 90\%} on every layer and above 95\% on lemma.
   \item All three parsers — \textbf{UDPipe, spaCy and Stanza} — produce consistent results across both treebanks.
   \item Stanza reaches \textbf{66 LAS and 97 UPOS} on the new data, a real step up for Uzbek parsing.
\end{itemize}

\noindent
The natural future work is to \textbf{extend} the treebank to conversational
and scholarly genres, to plug in \textbf{transformer-based parsers}, and to
explore \textbf{cross-lingual transfer} from Turkish, Kazakh, and Kyrgyz.

% ------------------------------------------------------------
%  SLIDE 16 — Thank you
% ------------------------------------------------------------
\section*{Slide 16 — Thank you}
\timing{13:45}{14:00}

That brings me to the end. I would like to thank \textbf{E.~Kuriyozov} for the
evaluation review, \textbf{K.~Madinabonu} and \textbf{M.~Shaydo} for the annotation
work, and \textbf{R.~Hulkaroy} for the morphological feature guidelines.

\say{The treebank is released in CoNLL-U format} — the URL is on the slide
and in the paper.

\pause

Thank you very much — I am happy to take your questions.

% ------------------------------------------------------------
%  Q & A prep — printed as a back page for the speaker only
% ------------------------------------------------------------
\newpage
\section*{Likely Q\&A — quick cues}

\noindent\textbf{Q. Why not just enlarge Uzbek-UT instead of building a new treebank?}\\
We considered it. The two efforts have \textbf{different annotation protocols}
and different licences. A separate, independent resource gives the
community a second reference point and makes \textbf{cross-checking and IAA
analysis possible}.

\vspace{0.4em}
\noindent\textbf{Q. The LAS gains are modest — only 3 to 4 points. Is that significant?}\\
Yes, with the same protocol and small test sets (200 / 100 sentences), a
\textbf{3--4 LAS gap is well outside seed variance}. We also see consistent
gains across \textbf{three independent parsers}, which is the stronger
evidence.

\vspace{0.4em}
\noindent\textbf{Q. Why no transformer baseline?}\\
We deliberately chose three \textbf{lightweight, reproducible} parsers so any
group can re-run the experiment. Transformer baselines are listed as
future work.

\vspace{0.4em}
\noindent\textbf{Q. How did you measure morphological feature IAA?}\\
We compute kappa and alpha \textbf{per token, per feature}, then aggregate.
The lower number (90.8\%) reflects ambiguity in derivational
suffixes, especially \textbf{verbal participles}.

\vspace{0.4em}
\noindent\textbf{Q. Is the treebank Cyrillic or Latin?}\\
\textbf{Latin script} only. All three source materials are in Latin.

\vspace{0.4em}
\noindent\textbf{Q. How are postpositions handled?}\\
As \textbf{separate tokens} attached with the \texttt{case} relation.
We follow the standard UD convention for Turkic languages.

\vspace{0.4em}
\noindent\textbf{Q. Will the data be on the official UD repository?}\\
Yes — we are preparing the submission for the next UD release cycle.

\end{document}
